% from ega4-10-fr.tex
\oldpage[IV-3]{116}

\begin{env}[10.10.5]
\label{IV.10.10.5-fr}
Enfin, on peut aussi considérer des ultra-préschémas sur un corps $k$
quelconque en conservant un langage qui reste proche de celui de Serre, et en
introduisant, comme chez Weil, une extension algébriquement close fixe $K$ de
$k$ (choisie assez grande, par exemple de degré de transcendance infini sur
$k$, pour avoir suffisamment de «~points génériques~» au sens de Weil). À tout
préschéma $X$ localement de type fini sur $k$, on associe alors l'ensemble
$S_K(X)=S(X\otimes_kK)$ des points de $X$ à valeurs dans $K$~; on a une
application canonique $j:S_K(X)\to X$ que l'on démontre être un
quasi-homéomorphisme quand on munit $S_K(X)$ de la topologie image réciproque
de celle de $X$ par $j$~; on munit $S_K(X)$ du faisceau de $k$-algèbres
$j^*(\mathcal O_X)$, et on obtient ainsi une sous-catégorie de la catégorie des
espaces $k$-annelés en anneaux locaux, que l'on pourrait appeler catégorie des
$(k,K)$-espaces préalgébriques. On peut montrer qu'on y peut encore définir des
produits et généraliser les résultats de (10.10.3) et (10.10.4)
($\mathcal A(X')$ étant ici remplacé par le faisceau $\mathcal A_K(X')$ des
germes d'applications de $X'$ dans $K$). Toutefois ce point de vue fait jouer
un rôle artificiel à un surcorps arbitrairement choisi de $k$, et nous ne le
signalons que pour le rejeter.
\end{env}

% (continued in ega4-11-fr.tex, before its first page marker)
\section{Propriétés topologiques des morphismes plats de présentation finie. Critères de platitude}
\label{section:IV.11-fr}

Alors que dans le \S~2 nous avons considéré les énoncés concernant la platitude
qui ne dépendent d'aucune hypothèse de finitude, et que le \S~6 étudie la notion
de platitude dans le cadre des préschémas localement noethériens (mais sans
hypothèse de finitude sur les morphismes), le présent paragraphe est consacré à
la notion de $f$-platitude dans le cas où le morphisme $f:X\to Y$ est
\emph{localement de présentation finie}. L'intérêt de la notion de morphisme
plat de présentation finie tient au fait que c'est elle qui semble exprimer
techniquement de la façon la plus adéquate la notion intuitive de «~famille de
préschémas algébriques paramétrée par un schéma $Y$~», dont l'étude est un des
objets principaux de la Géométrie algébrique. D'ailleurs, alors même qu'on ne
s'intéresserait au départ qu'au cas d'un schéma de base noethérien, il est
indispensable, pour certaines raisons techniques (par exemple pour certaines
applications de la théorie de la «~descente~», qui conduit à introduire des
schémas non nécessairement noethériens) de ne pas se limiter à ce cas, dès
qu'il s'agit de problèmes de nature essentiellement relative liés aux
morphismes localement de présentation finie. Nous suivrons systématiquement ce
principe, déjà étayé par les résultats des \S\S~8 et 9, dans toute la suite de ce
Chapitre, et même dans la suite de notre Traité, quitte à lui sacrifier à
l'occasion la simplicité de certaines démonstrations, que des hypothèses
noethériennes permettent parfois d'alléger\footnote{Ce principe s'inspire
également de la nécessité de donner droit de cité, comme «~espaces de
paramètres~» pour les familles de schémas algébriques, à des espaces annelés
(et même des «~topos~» annelés) quelconques, pour lesquels il ne peut plus être
question en général d'hypothèses noethériennes. Il semble assez clair que l'on
ne pourra plus éluder longtemps cette nouvelle extension de la Géométrie
algébrique, et il convient dès à présent de développer les notions et
techniques de nature «~relative~» de la théorie des schémas de sorte qu'elles
puissent s'adapter pratiquement telles quelles à ce cadre plus général.}. Dans
le présent paragraphe, cela nous conduit à reprendre, dans le contexte de la
«~présentation finie~» (notamment au n$^\circ$~3) certains énoncés de platitude,
déjà obtenus dans le contexte noethérien. L'outil technique essentiel pour
faire la réduction au cas noethérien est le théorème de compatibilité de la
platitude avec les limites projectives de préschémas (11.2.6), complétant les
résultats généraux du \S~8. Nous prouvons aussi en passant (11.3.1) un résultat
souvent utilisé

